\documentclass[10pt,a4paper]{amsart}
\usepackage[margin=19mm]{geometry}
\usepackage{amsmath,amssymb,mathtools}
\usepackage[scr=rsfs,cal=euler]{mathalfa}
\usepackage{xfrac}
\usepackage[unicode,hidelinks,bookmarks=false]{hyperref}
\newcommand{\ie}{i.e.\ }
\newcommand{\cf}{cf.\ }
\newcommand{\resp}{respectively\ }
\newcommand{\Subsection}{\subsection}
\providecommand{\nobreakdash}{\nobreak\hbox{-}}
\setlength{\emergencystretch}{2em}
\multlinegap0pt
\usepackage{bm}
\usepackage{bbm}
\usepackage{dsfont}
\usepackage{xspace}
\usepackage{cancel}
\usepackage{mathtools}
\usepackage{amsthm}
\makeatletter
\@ifundefined{lem}{\newtheorem{lem}{Lemma}}{}
\makeatother
\title[Parareal in time and spectral in space fast L1 quasilinear s]{Parareal in time and spectral in space fast L1 quasilinear subdiffusion solver}
\author{Josefa Caballero, {\L}ukasz P{\l}ociniczak, Kishin Sadarangani}
\date{Source: arXiv:2510.11023v1 (CC BY 4.0)}
\begin{document}
\maketitle

\noindent\emph{Source and licence.} Parareal in time and spectral in space fast L1 quasilinear subdiffusion solver. Version arXiv:2510.11023v1 (CC BY 4.0), distributed under the Creative Commons Attribution 4.0 International licence, as recorded in the arXiv licence metadata for this exact version and shown on the abstract page https://arxiv.org/abs/2510.11023v1. Full rights notice: licenses/arxiv-2510.11023-CC-BY-4.0.txt.\par\smallskip

\noindent\textit{Editorial scope.} Counted unit: the Lem whose source label is \texttt{lem:Gronwall} in the article (source0.tex lines 607--616, source label \texttt{lem:Gronwall}), delivered together with the article's own complete original proof of it (source0.tex lines 617--689, one \texttt{proof} environment of 73 lines). No mathematics is written, repaired or re-derived: statement and proof are the article's own text line for line. Required context retained uncounted: none. Every definition, display and label of those lines is retained unchanged. Omitted from this package and not redistributed: 1--606, 690--1375. Nothing inside a retained excerpt is omitted or altered: every retained body is delivered exactly as the article's lines are written, so no omission receipt of any kind accompanies this package. Citations: the retained excerpts carry no citation command at all, so no bibliography entry of the article is transcribed and no reference is redistributed. Reference adaptation: none was needed and none was made; every reference inside the delivered unit resolves inside it, so no unresolved reference or citation is produced. Printed slips kept literal: no printed word, symbol, hypothesis or number of the counted statement or of its proof was altered. Layout and class: the article's own class is \texttt{article} with packages including amssymb, amsmath, amsthm, algorithm, algpseudocode, fullpage, babel, graphicx, booktabs, xcolor, tikz, hyperref, autonum; the wrapper is the lane's pinned \texttt{amsart} editorial wrapper at 10pt on A4, which restates the theorem-like environments the retained text needs and adds the article's own macro definitions that the retained text needs (none), with extra wrapper packages (bm, bbm, dsfont, xspace, cancel, mathtools). The delivered numbering is the unit's own. Assets: no retained span contains a figure environment, a graphics-inclusion command, a PDF-page inclusion, a table or a TikZ picture; no image file is copied into this package, the package declares no asset dependency and no image file is redistributed with it. Licence: version arXiv:2510.11023v1 of this article is distributed under the Creative Commons Attribution 4.0 International licence, as recorded in the arXiv licence metadata for this exact version and shown on the abstract page https://arxiv.org/abs/2510.11023v1. A full rights notice is delivered at licenses/arxiv-2510.11023-CC-BY-4.0.txt. Characters: every retained excerpt is pure ASCII, so the delivered engine, XeLaTeX with its default Latin Modern text font, reports no missing character.
\begin{lem}\label{lem:Gronwall}
	Let $\left\{E^k_n\right\}_{n,k \in \mathbb{N}}$ be a positive sequence. If
	\begin{equation}
		E^{k+1}_{n+1} \leq a + b E^k_n + c E^{k+1}_{n}, \quad E^k_0 = 0, 
	\end{equation}
	for $a,b,c > 0$, then
	\begin{equation}\label{eqn:LemmaInequality}
		E^k_n \leq a \sum_{j=0}^{k-1} \sum_{i=0}^{n-j-1} \binom{i+j}{j} c^i b^j + E^0_n b^k \sum_{j=k}^{n-1} \binom{j-1}{k-1} c^{j-k}.
	\end{equation}
\end{lem}

\begin{proof}
	The proof proceeds in three steps:
	\begin{enumerate}
		\item Formulate auxiliary \emph{equation}.
		\item Solve the homogeneous equation.
		\item Solve the non-homogeneous problem.
	\end{enumerate}
	\medskip\noindent\textbf{Ad. 1.} First, define $\{f^k_n\}_{n,k\in\mathbb{N}}$ to be the solution of the following equation
	\begin{equation}
		f^{k+1}_{n+1} = a + b f^k_n + c f^{k+1}_{n}, \quad f^k_0 = 0, \quad f^0_n = E^0_n.
	\end{equation}
	We will show by induction that $E^k_n \leq f^k_n$. This shows that it is sufficient only to solve the equation. By the definition of the initial values $f_0^k$ and $f^0_n$ we see that the first step in the induction is satisfied. Now, assume that $E^i_j \leq f^i_j$ for all $0\leq i \leq k+1$ and $0\leq n \leq n$. By this assumption and the definition of $E^{k+1}_{n+1}$, we have
	\begin{equation}
		E^{k+1}_{n+1} \leq a + b E^k_n + c E^{k+1}_n \leq a + b f^{k}_n + c f^{k+1}_n = f^{k+1}_{n+1}, 
	\end{equation}
	and the induction is complete. 
	
	\medskip\noindent\textbf{Ad. 2.} We proceed to solving the main equation. By linearity, we can first consider the homogeneous version
	\begin{equation}
		f^{k+1}_{n+1} = b f^k_n + c f^{k+1}_{n}, \quad f^k_0 = 0, \quad f^0_n = E^0_n.
	\end{equation}
	We will that
	\begin{equation}
		f^k_n = E^0_n b^k \sum_{j=k}^{n-1} \binom{j-1}{k-1} c^{j-k}.
	\end{equation}
	This formula can be found by several techniques: "guessing" based on numerical computations, using generating functions, or combinatorial counting argument. To verify the above, we have to check whether the ansatz satisfies the equation. The "initial condition" is 
	\begin{equation}
		f^0_n = E^0_n \binom{-1}{-1} = E^0_n,
	\end{equation}
	since all other binomial coefficients vanish. On the other hand, the right-hand side of the homogeneous equation is
	\begin{equation}
		\begin{split}
			b f^k_n + c f^{k+1}_{n} 
			&= E^0_n \left( b^{k+1} \sum_{j=k}^{n-1} \binom{j-1}{k-1} c^{j-k} + b^{k+1} \sum_{j=k+1}^{n-1} \binom{j-1}{k} c^{j-k-1+1}\right) \\
			&= E^0_n b^{k+1} \left( \sum_{j=k}^{n-1} \binom{j-1}{k-1} c^{j-k} + \sum_{j=k+1}^{n-1} \binom{j-1}{k} c^{j-k}\right).
		\end{split}
	\end{equation} 
	Now, we separate the first term in the first sum and merge
	\begin{equation}
		b f^k_n + c f^{k+1}_{n} = E^0_n b^{k+1} \left(1+ \sum_{j=k+1}^{n-1} \left[\binom{j-1}{k-1} + \binom{j-1}{k}\right]c^{j-k}\right) = E^0_n b^{k+1} \left(1+ \sum_{j=k+1}^{n-1}\binom{j}{k} c^{j-k}\right),
	\end{equation} 
	where we used the fundamental recurrence identify for the binomial coefficient. Finally, we change the summation index $i = j+1$ and obtain
	\begin{equation}
		b f^k_n + c f^{k+1}_{n} = E^0_n b^{k+1} \left(1+ \sum_{i=k+2}^{n-1}\binom{i-1}{k} c^{i-k-1}\right) = E^0_n b^{k+1} \sum_{i=k+1}^{n-1}\binom{i-1}{k} c^{i-k-1} = f^{k+1}_{n+1},
	\end{equation} 
	which concludes the proof of this part. 
	
	\medskip\noindent\textbf{Ad. 3.} Now, we turn to the nonhomogeneous problem with vanishing initial condition
	\begin{equation}
		f^{k+1}_{n+1} = a + b f^k_n + c f^{k+1}_{n}, \quad f^k_0 = 0, \quad f^0_n = 0.
	\end{equation}
	We claim that
	\begin{equation}
		f^k_n = a \sum_{j=0}^{k-1} \sum_{i=0}^{n-j-1} \binom{i+j}{j} c^i b^j.
	\end{equation}
	First, we see that by the convention in which the upper index of the summation is smaller than the lower, we have $f^0_n = 0$ and $f^k_0 = 0$. Moreover, the ansatz satisfies the equation as can be seen by changing the summation indices
	\begin{equation}
		\begin{split}
			a + b f^k_n + c f^{k+1}_{n} 
			&= a\left(1+\sum_{j=0}^{k-1} \sum_{i=0}^{n-j} \binom{i+j}{j} c^i b^{j+1} + \sum_{j=0}^{k} \sum_{i=0}^{n-j} \binom{i+j}{j} c^{i+1} b^j\right) \\
			&= a\left(1+\sum_{j=1}^{k} \sum_{i=0}^{n-j} \binom{i+j-1}{j-1} c^i b^{j} + \sum_{j=0}^{k} \sum_{i=1}^{n-j} \binom{i+j-1}{j} c^{i} b^j\right).
		\end{split}
	\end{equation}
	Then, by noticing that since $\binom{p}{q} = 0$ if $q > p$, all the $(i,j)$-th terms of the above sums vanish apart the one with $i = j = 0$ with a value $1$, we can artificially add and subtract it
	\begin{equation}
		\begin{split}
			a + b f^k_n + c f^{k+1}_{n} 
			&= a\left(1-1+\sum_{j=0}^{k} \sum_{i=0}^{n-j} \binom{i+j-1}{j-1} c^i b^{j} + \sum_{j=0}^{k} \sum_{i=0}^{n-j} \binom{i+j-1}{j} c^{i} b^j\right) \\
			&= a\sum_{j=0}^{k} \sum_{i=0}^{n-j} \left[\binom{i+j-1}{j-1} + \binom{i+j-1}{j}  \right] c^i b^{j} = a\sum_{j=0}^{k} \sum_{i=0}^{n-j} \binom{i+j}{j} c^i b^{j} = f^{k+1}_{n+1},
		\end{split}
	\end{equation}
	which completes the proof of the lemma.
\end{proof}

\end{document}
